\section{Reduction to a supported nef boundary}\label{sec:reduction}

The positive-Iitaka-dimensional case is an application of a known
K\"ahler abundance theorem.  The purpose of this section is to prepare
the remaining case for the two boundary arguments: from a nonzero divisor
of a section in Iitaka dimension zero, we construct a nef dlt fourfold
whose adjoint has a nonzero effective representative supported on exactly
its reduced floor.  We also construct the particular log resolution used
to index and compare all strata of that floor.

\begin{proposition}[Supported nef model]\label{prop:supported-model}
Under the hypotheses of Theorem~\ref{thm:main}, suppose
\(\kappa(X,D)=0\) and the normalized rational divisor \(M\) of a nonzero plurisection of
\(D\) is nonzero.  Then there are a normal ordinary \(\Q\)-factorial
compact K\"ahler dlt fourfold \((V,B)\), with effective rational boundary,
and a nonzero effective rational \(\Q\)-Cartier divisor \(P\) such that
\[
 A=K_V+B\text{ is analytically nef},\qquad
 A\sim_\Q P,\qquad \Supp P=\Supp\floor B,\qquad \kappa(V,A)=0.
\]
The pair has the projective resolution described in
Lemma~\ref{lem:special-resolution}.
\end{proposition}

The resolution has globally smooth distinct SNC strict boundary and
exceptional components, all exceptional crepant coefficients are below
one, and it is generically an isomorphism on the image of every
intersection component of distinct strict floor primes.  These properties
let the floor argument index its strata by actual boundary intersections.
The proposition will follow from the construction below; the final step
uses the original nef adjoint to show that \(P\) cannot disappear.

We use the following analytic negativity theorem at several points.
If \(h:W\to Z\) is a proper bimeromorphic morphism of normal irreducible
analytic spaces and \(E\) is a rational \(\Q\)-Cartier divisor with \(-E\)
relatively nef, then
\[
 E\geq0\quad\Longleftrightarrow\quad h_*E\geq0.
\]
This is \cite[Lemma~1.3]{Wang2021}, after clearing an index.  In particular,
an exceptional relatively nef divisor is nonpositive.  When \(h\) is
projective, nonnegative degrees on its contracted curves give the relative
nefness used in this theorem; see \cite[Appendix~B]{Wang2021}.  This relative
curve test will not be used as a definition of nefness on a nonprojective
compact K\"ahler space.

\subsection{Positive Iitaka dimension}

\begin{proposition}\label{prop:positive-kappa}
Under the hypotheses of Theorem~\ref{thm:main}, if \(\kappa(X,D)\geq1\),
then \(D\) is semiample on \(X\).
\end{proposition}
\begin{proof}
Das--Hacon--P\u{a}un's dlt modification theorem
\cite[Theorem~6.1]{DasHaconPaun2024} applies to a compact K\"ahler lc
fourfold with effective rational boundary and \(\Q\)-Cartier adjoint.
It gives a projective bimeromorphic morphism
\[
 g:(X',\Delta')\longrightarrow(X,\Delta)
\]
with \(X'\) compact K\"ahler and ordinary \(\Q\)-factorial, the pair
\((X',\Delta')\) dlt, and the adjoint crepant.  The input is klt, so equality
of discrepancies makes the output klt as well.  The crepant equality is
an equality of actual rational lines: pull a local pluriadjoint frame to
a common resolution with smooth source as in \eqref{eq:crepant}, compare the crepant
coefficients, and extend the identical meromorphic map across codimension
two on the normal \(X'\).  Thus for a common index
\[
 \OO_{X'}(m(K_{X'}+\Delta'))\simeq g^*\OO_X(mD).
\]
There is no possible undetected flat-line difference in this comparison.

The pulled-back line is analytically nef by the metric criterion.  Also,
normality and projection formula identify all its sections with those of
\(\OO_X(mD)\).  The meromorphic maps agree on the common dense open, so
their Iitaka dimensions agree.  Theorem~4.1 of
H\"oring--Lazi\'c--Lehn \cite{HoeringLazicLehn2025} states that a nef
adjoint of a compact ordinary \(\Q\)-factorial K\"ahler klt pair of
positive Iitaka dimension is semiample, unconditionally in dimension at
most four.  Its analytic positivity and canonical-sheaf conventions are
the ones in use here.  Apply it to \((X',\Delta')\), then enlarge the
generated degree to a multiple of the original index.  All sections are
pullbacks.  A section nonzero at a chosen point above \(x\in X\) descends
to a section nonzero at \(x\); Nakayama's lemma gives the evaluation
surjectivity on \(X\).
\end{proof}

\subsection{A log-smooth supported representative}

We next begin with a divisor of a section.  No nefness is needed for the
preparation in this subsection.  The later proof that the support survives
the minimal model program will use the nefness of the original adjoint.

\begin{lemma}\label{lem:supported-preparation}
Let \((X,\Delta)\) be a normal connected compact K\"ahler klt pair with
effective rational boundary and actual \(\Q\)-Cartier adjoint \(D\).
Assume \(\kappa(X,D)=0\), and let \(M\geq0\) be a rational
\(\Q\)-Cartier divisor with \(M\sim_\Q D\).  There are a projective
modification \(p:Y\to X\) with \(Y\) smooth compact K\"ahler, an effective rational
SNC boundary \(B_Y\), and an effective rational divisor \(P_Y\) such that
\[
 K_Y+B_Y\sim_\Q P_Y,\qquad
 \Supp P_Y=\Supp\floor{B_Y},\qquad
 \kappa(Y,K_Y+B_Y)=0.
\]
The SNC support has globally smooth distinct components, and \(B_Y\) has
coefficient one on every \(p\)-exceptional prime and on every strict
transform of a prime in \(M\).
\end{lemma}
\begin{proof}
Principalize the reduced coherent ideal of \(\Supp(\Delta+M)\), using
\cite[Theorem~2.13 and Remark~2.14]{DasHaconPaun2024}.  This use of an
ideal is important: the original boundary need not be \(\Q\)-Cartier.
The resulting modification has smooth source, is projective, and has locally normal
crossing strict and exceptional support.  Mark each of its finitely many
global prime components separately as an ordered Cartier boundary and
apply the ordered-boundary resolution of
\cite[Theorems~1.1.3 and~1.1.13]{Temkin2017}.  Its final indexed
components and their intersections are smooth, and its complete support
consists of strict transforms and new exceptional components
\cite[Lemmas~2.1.10 and~2.2.9(iv)]{Temkin2017}.  Denote the composite by
\(p:Y\to X\).  It is projective over the compact base.  A relative ample
metric plus a sufficiently large pullback of a K\"ahler form makes \(Y\)
compact K\"ahler.

Let \(G_Y\) be the crepant subboundary for \(D\), defined by the actual
meromorphic pluriadjoint pullback.  Then \(p_*G_Y=\Delta\), and klt gives
\(\operatorname{coeff}_E(G_Y)<1\) at every prime.  Define \(B_Y\) only
after the preceding resolution: assign coefficient one to the actual
\(p\)-exceptional primes and the strict transforms of primes in \(M\),
and retain the coefficients of \(\Delta\) on the other strict boundary
primes.  An exceptional label in the resolution bookkeeping that is not
an actual exceptional divisor does not change this rule.  This is an
effective rational SNC boundary, hence dlt.

Set
\begin{equation}\label{eq:prepared-P}
 P_Y=p^*M+(B_Y-G_Y).
\end{equation}
Pullback of the holomorphic equation of a Cartier multiple of \(M\)
makes \(p^*M\) effective.  On a strict prime of \(M\), the second summand
has coefficient \(1-\operatorname{coeff}(\Delta)>0\); on an exceptional
prime it has coefficient \(1-\operatorname{coeff}(G_Y)>0\); on any other
strict boundary prime it is zero.  Thus \(P_Y\) is effective and its
reduced support is exactly the floor of \(B_Y\).  The actual identity
\(p^*D\sim_\Q p^*M\), together with \eqref{eq:crepant}, gives
\(K_Y+B_Y\sim_\Q P_Y\).

Choose an integer \(t>0\) at least the ratios of the coefficients of
\(P_Y\) to those of \(p^*M\) on every nonexceptional prime of \(M\).
The positive part of \(P_Y-tp^*M\) is then an effective rational
exceptional divisor \(E\), and
\[
 P_Y\leq tp^*M+E.
\]
For every sufficiently divisible \(n\), exceptional Hartogs extension
and projection formula give
\[
 H^0\bigl(Y,\OO_Y(n(tp^*M+E))\bigr)
 =H^0\bigl(X,\OO_X(ntM)\bigr).
\]
The space on the right has dimension one: it is nonzero and
\(\kappa(X,D)=0\).  The displayed divisor inequality bounds
\(h^0(Y,\OO_Y(nP_Y))\) by one, while effectivity makes it nonzero.
This controls every sufficiently divisible degree.  If two independent
sections existed in another Cartier degree, their powers \(s_0^v\) and
\(s_0^{v-1}s_1\) would remain independent in a degree divisible by the
fixed common index, a contradiction.  Therefore \(\kappa(Y,P_Y)=0\).
\end{proof}

\subsection{Projective steps with K\"ahler targets}

Apply the preceding lemma in dimension four.  We follow the supported
construction of Das--Hacon--P\u{a}un
\cite[Theorem~7.2]{DasHaconPaun2024} from the compact ordinary
\(\Q\)-factorial K\"ahler dlt pair \((Y,B_Y)\).  We isolate its
three-dimensional floor-contraction input through
Subsection~\ref{subsec:actual-floor-split} and use
Proposition~\ref{prop:prime-floor-extension} for the ambient contraction.
The construction below gives projective steps with K\"ahler targets and
supplies the graphs needed for the discrepancy comparison.

Write \((X_i,B_i)\) for a nonnef running pair and
\(A_i=K_{X_i}+B_i\).  At the start of this step, induction from
\(P_Y\) gives an effective rational divisor \(P_i\) with
\begin{equation}\label{eq:running-supported-representative}
 P_i\sim_\Q A_i,\qquad
 \Supp P_i=\Supp\floor{B_i}.
\end{equation}
The equivalence is an isomorphism of actual rational holomorphic lines.
It holds initially by the supported-model construction, and the
one-step transform identity proved below supplies it at the next step
only after the current step has been completed.  Thus the representative
used now is already available before constructing the current target.

The floor-contraction input in the proof of
Theorem~7.2, using its Corollary~5.6, is the following assertion.
There are a global irreducible floor prime \(T\), a K\"ahler form
\(\omega_i\) on \(X_i\), a projective contraction
\(\varphi:T\to W\) to a normal compact K\"ahler space with
\(\varphi_*\OO_T=\OO_W\), and a K\"ahler form \(\omega_W\).
The class \(\alpha_i=c_1(A_i)+[\omega_i]\) is nef and big but not
K\"ahler, and
\begin{equation}\label{eq:floor-contraction-class}
 c_1(A_i|_T)+[\omega_i|_T]=\varphi^*[\omega_W].
\end{equation}
Here the equality is in the local-potential Bott--Chern group.  The
contracted compact curves are exactly those whose classes on \(T\) lie in
\[
 F_T=\bigl(c_1(A_i|_T)+[\omega_i|_T]\bigr)^\perp
       \cap\operatorname{NA}(T).
\]
This face is generated by finitely many curve classes whose images in
\(X_i\) lie in one \(A_i\)-negative ray
\(R_i\subset\alpha_i^\perp\), with \(T\cdot R_i<0\).
For the global form of this selection, apply Lemma~7.1 in the setup of
Theorem~7.2 together with Claim~7.3 and its proof.  We choose one
representative per proportionality ray among the countably many compact
curve classes to meet the lemma's nonproportionality hypothesis.  A
K\"ahler degree is positive on each curve, so proportional effective
representatives differ by a positive scalar.  Lemma~7.1 permits at most
one representative on which the selected class vanishes, and Claim~7.3
and its proof supply it.  Consequently the selection gives the global
implication
\begin{equation}\label{eq:global-support-zero-ray}
 \alpha_i\cdot C=0\quad\Longrightarrow\quad [C]\in R_i
 \quad\text{for every compact irreducible curve }C\subset X_i.
\end{equation}
For a positive global Cartier multiple
\(mT\), the line \(\OO_T(-mT)\) is \(\varphi\)-ample.
The use of Corollary~5.6 here is the actual-line specialization proved
in Subsection~\ref{subsec:actual-floor-split}, with the inputs stated in
Subsection~\ref{subsec:threefold-contraction-bridge}.
In particular, the restriction to \(T\) need not be big and \(T\)
need not be \(\Q\)-factorial.  The numerical selection of \(T\),
the ray, and the negative normal line is the one in the supported
construction.  Its threefold contraction is supplied by the bridge,
using Proposition~\ref{prop:prime-floor-extension} and the external
results stated there.

Write \(B_i=T+C\), where \(C\geq0\) has no component \(T\).
Ordinary \(\Q\)-factoriality makes \(C\) a global rational
\(\Q\)-Cartier divisor.  For a positive rational
\(\varepsilon<1\), put
\[
 B_i(\varepsilon)=T+(1-\varepsilon)C,\qquad
 A_i(\varepsilon)=A_i-\varepsilon C
                 =K_{X_i}+B_i(\varepsilon).
\]
The second identity is an identity of actual rational adjoint lines.
The perturbed pair is plt.  Indeed, choose a log resolution
\(f:U\to X_i\) with smooth source that witnesses the standard dlt
criterion for \((X_i,B_i)\).
Its strict boundary and exceptional primes have simple normal crossings,
and all exceptional crepant coefficients are below one.  If \(G\) is its
crepant boundary, the new boundary on this resolution is
\(G-\varepsilon f^*C\).  The divisor \(f^*C\) is effective: pull back
a local holomorphic equation for an effective Cartier multiple of \(C\).
It follows that all exceptional coefficients remain below one.  The
strict transform of \(T\) is the only coefficient-one prime, and all
other strict coefficients are below one.  The plt criterion on this
resolution now applies.  This is the ordinary Koll\'ar--Mori convention
used in Section~2 of \cite{DasHacon2024OnMMP}.  If \(C=0\), the same
argument applies to the unchanged pair \((X_i,T)\).

Choose a smooth Hermitian metric on an actual Cartier multiple of \(C\),
and let \(\theta_C\) be its normalized curvature form.  On a finite
relatively compact local embedding cover of the compact \(X_i\), choose
smooth ambient extensions of the local potentials and metric weights.
After shrinking the charts, the Levi forms for \(\omega_i\) have a
uniform positive lower bound on their compact closures, while those for
\(\theta_C\) are bounded.  Thus \(\omega_i+\varepsilon\theta_C\) is K\"ahler for
one sufficiently small positive rational \(\varepsilon\), chosen for
this running step.  No positivity of \(C|_T\) is asserted.  Equation
\eqref{eq:floor-contraction-class} gives
\begin{equation}\label{eq:plt-contraction-class}
 -c_1(A_i(\varepsilon)|_T)+\varphi^*[\omega_W]
       =[\omega_i|_T+\varepsilon\theta_C|_T].
\end{equation}
The right side is K\"ahler.  Locally on \(W\), a potential for
\(\omega_W\) can be pulled back, so this equality says precisely that
\(-A_i(\varepsilon)|_T\) is relatively K\"ahler over \(W\).
After clearing its index, the same local metric weights restrict to
positive weights on every full \(\varphi\)-fiber.  The proper fiberwise
criterion \cite[Corollary~1.12, Definition~3.1 and Remark~3.2]{Fujino2026Fujiki}
therefore makes this actual rational line \(\varphi\)-ample.  The other
required positivity, that of \(-T|_T\), is unchanged.  With the boundary
\((1-\varepsilon)C\), these are the hypotheses of
Proposition~\ref{prop:conormal-layers}; its conormal direct-image
vanishings hold for the present floor contraction.

Proposition~\ref{prop:prime-floor-extension} now constructs a proper
bimeromorphic morphism \(c:X_i\to Z_i\) to a normal compact analytic
space.  It restricts to \(\varphi\) on \(T\), embeds \(W\) as the
reduced image of \(T\), has exactly the \(\varphi\)-fiber sets over
\(W\), and is an isomorphism away from \(T\).  Its fibers are
connected.  A sufficiently divisible global line \(\OO_{X_i}(-\ell T)\)
is \(c\)-ample, and the actual rational line
\(-A_i(\varepsilon)\) is \(c\)-ample as well.  The target is in
Fujiki's class \(\mathcal C\).  The proposition proves these assertions
through the analytic thickenings and the full-fiber criterion.  For this
constructed contraction, we use the rationality assertion of
\cite[Remark~5.11]{DasHacon2024OnMMP}.  The target K\"ahler class is
proved below.

Every curve contracted by \(c\) is a \(\varphi\)-contracted curve in
\(T\), and therefore has class in \(R_i\).  Conversely a curve in
\(R_i\) lies in \(T\), since \(T\cdot R_i<0\) whereas the effective
Cartier section of a multiple of \(T\) has nonnegative degree on the
normalization of any curve not contained in \(T\).  The floor assertion
then contracts it.  The constructed morphism thus contracts precisely
the compact curves whose classes lie in the original negative ray.
The boundary transported in the supported
program remains \(B_i\); the plt boundary is used only to construct \(c\).

\begin{claim}[K\"ahler positivity on the current target]
\label{claim:outer-target-positivity}
Retain the current pair \((X_i,B_i)\), its effective rational divisor
\(P_i\) with the actual line identity in
\eqref{eq:running-supported-representative}, the nef and big class
\(\alpha_i=c_1(A_i)+[\omega_i]\) with \(\omega_i\) K\"ahler, and the
selected ray and floor in
\eqref{eq:floor-contraction-class}--\eqref{eq:global-support-zero-ray}.
Let \(c:X_i\to Z_i\) be the projective contraction with connected fibers
just constructed.  It is an isomorphism away from \(T\), has the embedded
floor image \(W\) and the exact \(\varphi\)-fiber sets, and contracts
precisely the curves in \(R_i\).  Its target is a normal compact analytic
space in class \(\mathcal C\) with rational singularities.  For every other floor prime,
retain the projective contraction with connected fibers supplied by
Subsection~\ref{subsec:actual-floor-split}, whose target is normal compact
K\"ahler and whose pulled-back K\"ahler class is the restricted
\(\alpha_i\).
Then there is a K\"ahler form \(\omega_{Z_i}\) on \(Z_i\) such that
\begin{equation}\label{eq:outer-target-kahler}
 \alpha_i=c^*[\omega_{Z_i}]
 \quad\text{in }H^{1,1}_{BC}(X_i),
\end{equation}
where the group is defined by local smooth potentials.
\end{claim}
\begin{proof}
Bott--Chern descent first supplies a target class whose pullback is
\(\alpha_i\).  We prove that this class is nef, contains a current
dominating a positive Hermitian form, and has positive top intersection
on every positive-dimensional compact reduced subspace.  These are the
three conditions in \cite[Theorem~2.29]{DasHaconPaun2024} that make the
target class K\"ahler.
\paragraph{The descended class and its floor restriction.}
In this paragraph and the next three, write \(X=X_i\), \(Z=Z_i\), and
\(\alpha=\alpha_i\).  We use
\(dd^c=\frac{i}{2\pi}\partial\bar\partial\), the normalization for
which \(dd^c\phi\) is the normalized Chern curvature of a metric
locally written as \(e^{-\phi}\).  A fixed positive rescaling of potentials gives the
\(i\partial\bar\partial\) convention in the cited analytic inequalities.

Both ends of \(c\) are normal compact spaces in class \(\mathcal C\)
with rational singularities, and \(\alpha\) is zero on every contracted
curve.  The Bott--Chern descent
\cite[Lemma~2.11]{DasHacon2024} therefore gives a class \(\beta\) on
\(Z\) and a smooth real closed representative \(\beta_0\), with local
smooth potentials, such that
\begin{equation}\label{eq:outer-descended-class}
 c^*\beta=\alpha\quad\text{in }H^{1,1}_{BC}(X).
\end{equation}
Here and below the group is the local-potential Bott--Chern group.
Its smooth representative description, including adjustment by a global
smooth \(dd^c\)-potential, is
\cite[Definition~3.1 and Remark~3.2]{HoeringPeternell2016}.

Let \(i_W:W\hookrightarrow Z\) be the embedded floor image and put
\(\delta=i_W^*\beta\).  The actual identity
\(c|_T=i_W\circ\varphi\), \eqref{eq:floor-contraction-class}, and
\eqref{eq:outer-descended-class} give
\begin{equation}\label{eq:outer-floor-image-class}
 \varphi^*\delta=\alpha|_T=\varphi^*[\omega_W]
 \quad\text{in }H^{1,1}_{BC}(T).
\end{equation}
Thus the pullback of \(\delta-[\omega_W]\) is zero and hence nef.
Apply \cite[Lemma~2.38]{DasHaconPaun2024} only to
\(\varphi:T\to W\): it is proper and surjective, and both spaces are
normal compact K\"ahler.  The lemma makes \(\delta-[\omega_W]\) nef.
For a smooth representative \(e\) of that difference, choose a nef
approximation \(e+dd^c u\geq-\omega_W/2\).  Then
\(\omega_W+e+dd^c u\) is a K\"ahler representative of \(\delta\).
In particular its restriction gives a positive current in the class
\(\beta|_V\) and a positive top integral for every positive-dimensional
irreducible \(V\subset W\).  Only the nefness of
\(\delta-[\omega_W]\) enters this deduction.

\paragraph{Nefness on the target.}
For every positive-dimensional irreducible reduced compact subspace
\(V\subset Z\), we first produce a current
\begin{equation}\label{eq:outer-restriction-current}
 T_V=\beta_0|_V+dd^c q_V\geq0
\end{equation}
with a global real distribution \(q_V\) on \(V\).  Currents and their
positivity on the pure-dimensional reduced \(V\) are defined by duality
with smooth test forms from local ambient embeddings; see
\cite[Section~1, Definitions~1.1--1.2, pp.~14--15]{Demailly1985}.
We will use exactly this positive-current meaning of pseudoeffectivity
on a possibly nonnormal subspace.  The preceding K\"ahler representative
of \(\delta\) gives \eqref{eq:outer-restriction-current} for
\(V\subset W\), with a smooth potential.

Suppose \(V\not\subset W\), and let \(V'\subset X\) be the closure
of the inverse image of \(V\setminus W\).  The isomorphism
\(X\setminus T\simeq Z\setminus W\) makes \(V'\) irreducible and
its map to \(V\) proper and bimeromorphic.  Choose a projective
resolution \(r:U\to V'\) with smooth source.  The restricted K\"ahler
form makes \(V'\) a compact K\"ahler space.  A relatively positive
metric for the projective resolution, plus a sufficiently large pullback
of this form, makes \(U\) a compact K\"ahler manifold.  This metric
construction applies to the possibly nonnormal \(V'\).

Let \(\nu:V^\nu\to V\) be normalization.  The dominant map from the
normal \(U\) factors through \(\nu\), giving a proper bimeromorphic
map \(p:U\to V^\nu\), hence a modification between normal compact
spaces.  For the other map \(j:U\to X\), functoriality of the actual
maps and \eqref{eq:outer-descended-class} gives
\[
 p^*\nu^*(\beta|_V)=j^*\alpha.
\]
The class on the right is nef.  Indeed, pull back the smooth nef
approximations from \(X\); on compact \(U\), the pullback of their
reference form is bounded above by a fixed multiple of a K\"ahler form
\(\eta\) on \(U\), so rescaling the error gives the nef inequalities.

Set \(\gamma=p^*\nu^*(\beta_0|_V)\).  Choose smooth approximations
\(\gamma+dd^c f_n\geq-n^{-1}\eta\).  The closed positive currents
\(\gamma+dd^c f_n+n^{-1}\eta\) have bounded \(\eta\)-mass by
Stokes' theorem.  Weak compactness gives a positive closed limit in
\([\gamma]\).  The \(\partial\bar\partial\)-lemma on the compact
K\"ahler manifold \(U\), and the global potential description of a
positive current, give a global quasi-plurisubharmonic \(q_U\) with
\[
 \gamma+dd^c q_U\geq0.
\]
These are also the closed-cone statement following Definition~1.6 and
the potential description in (3.1) and its following paragraph in
\cite[pp.~1253 and 1260]{DemaillyPaun2004}.  Apply
\cite[Corollary~2.32 and its proof, p.~18]{DasHaconPaun2024} to \(p\), with target form
\(\nu^*(\beta_0|_V)\) and constant zero in its current inequality.
The locally integrable, locally bounded above \(q_U\) supplies the
required potential.  We obtain a global quasi-plurisubharmonic
\(q_\nu\) satisfying
\begin{equation}\label{eq:outer-normalized-current}
 T_\nu=\nu^*(\beta_0|_V)+dd^c q_\nu\geq0.
\end{equation}
This applies the normal-modification descent to the displayed
local-potential inequality on the still general normal compact target.

It remains to push through the finite normalization.  Locally write
\(\beta_0|_V=dd^c h\) with \(h\) smooth.  The function
\(\nu^*h+q_\nu\) is locally integrable and locally bounded above, and
its Hessian is \(T_\nu\geq0\).  On normal \(V^\nu\), its upper
regularization is plurisubharmonic.  The finite trace is weakly
plurisubharmonic and its Hessian is the positive current \(\nu_*T_\nu\)
by \cite[Corollary~1.11 and Proposition~1.13(a), p.~22]{Demailly1985}.
Since \(\nu\) has generic degree one, change of variables off the
proper analytic exceptional sets gives
\(\nu_*\nu^*(\beta_0|_V)=\beta_0|_V\) as currents; those sets have
measure zero for the smooth top-degree integrands.  Proper push of
distributions commutes with \(dd^c\).  Hence
\begin{equation}\label{eq:outer-finite-trace}
 T_V=\nu_*T_\nu
     =\beta_0|_V+dd^c(\nu_*q_\nu)\geq0.
\end{equation}
The traced distributions are locally integrable and glue because
\(q_\nu\) is global.  On a locally reducible \(V\), these local
potentials may be only weakly plurisubharmonic; the argument below uses
their associated positive currents.

We now apply the sufficient direction of the restriction criterion
\cite[Theorem~2.36 and Remark~2.37]{DasHaconPaun2024}, spelling out its
singular-current input.  Proposition~3.3(iv), p.~1262, and the following
paragraph, pp.~1262--1263, of \cite{DemaillyPaun2004}
applies to a compact complex space and a smooth class containing a
closed positive current.  It gives nefness if the restrictions to the
irreducible components of every positive Lelong level set are nef.
Induct on the dimension of each irreducible \(V\subset Z\); points are
automatic.  For a positive-dimensional \(V\), use
\eqref{eq:outer-restriction-current}.  The Siu analyticity assertion
accompanying that proposition makes each positive Lelong level set
analytic, and it is proper.  Indeed, if a fixed positive level filled \(V\), pack
disjoint radius-\(r\) balls in a coordinate ball of \(V_{\rm reg}\).
The Lelong mass lower bound on each ball would force the locally finite
mass of the \((1,1)\)-current to grow at least as a positive multiple
of \(r^{-2}\) when \(r\) tends to zero.  This is impossible.  The level
set components therefore have smaller dimension, and their restricted
classes are nef by induction.  The proposition makes \(\beta|_V\) nef.
Taking \(V=Z\) proves analytic nefness of \(\beta\).  This uses the
positive-current formulation above also on nonnormal subspaces, where
the local potentials may be only weakly plurisubharmonic.

\paragraph{Bigness from the supported section.}
Use the effective actual representative already present in
\eqref{eq:running-supported-representative}.  Choose \(m>0\) clearing
the indices and an isomorphism of holomorphic lines
\[
 L=\OO_X(mA_i)\simeq\OO_X(mP_i).
\]
The canonical section of the right side gives a nonzero holomorphic
section \(s\) of \(L\) with divisor \(mP_i\).  Choose a smooth metric
\(h\) on \(L\).  In a local frame write \(h=e^{-\phi}\), \(s=f\),
and put
\[
 \vartheta=\frac1m dd^c\phi,\qquad
 \ell=\frac1m\log|s|_h^2.
\]
The first formula defines the normalized global curvature form.
The second is a global quasi-plurisubharmonic function, locally
\(m^{-1}(\log|f|^2-\phi)\).  The holomorphic \(f\) extends in a local
ambient embedding and is not identically zero on any nonempty open
subset of that chart in \(X\).  Local integrability is Proposition~1.8
of \cite[p.~19]{Demailly1985}, and
the unnumbered prose after its proof gives the current positivity
\[
 \vartheta+dd^c\ell=\frac1m dd^c\log|f|^2\geq0.
\]
The smooth form \(\vartheta+\omega_i\) represents
\(\alpha=c^*\beta\), so the smooth local-potential representative
description supplies a global smooth \(v\) with
\(\vartheta+\omega_i=c^*\beta_0+dd^c v\).  Thus
\begin{equation}\label{eq:outer-supported-current}
 c^*\beta_0+dd^c(v+\ell)\geq\omega_i.
\end{equation}
This is a global quasi-plurisubharmonic potential obtained from the
supported section at this running step.

Fix any smooth positive Hermitian form \(g_Z\) on \(Z\) in local
ambient embeddings.  Local holomorphic coordinate lifts of \(c\) make
\(c^*g_Z\) semipositive in those embeddings.  A finite relatively
compact cover of \(X\) and comparison with the positive definite
ambient representatives of \(\omega_i\) give one \(C>0\) with
\(c^*g_Z\leq C\omega_i\).  Corollary~2.32 of
\cite{DasHaconPaun2024} applies to the normal modification \(c\) and
\eqref{eq:outer-supported-current}.  It gives a global
quasi-plurisubharmonic, hence \(L^1\), function \(v_Z\) with
\begin{equation}\label{eq:outer-current-big}
 \beta_0+dd^c v_Z\geq C^{-1}g_Z.
\end{equation}
This is the required dominating current in the selected target class.

\paragraph{Top intersections and the other floor.}
Let \(V\subset Z\) be irreducible of dimension \(k>0\).  The case
\(V\subset W\) was proved using the K\"ahler class \(\delta\).
Otherwise use its strict transform \(V'\) above.  The proper integration
cycle identity \(c_*[V']=[V]\), projection formula, and
\eqref{eq:outer-descended-class} give
\begin{equation}\label{eq:outer-strict-transform}
 \int_{V_{\rm reg}}\beta_0^k
   =\int_{V'_{\rm reg}}\alpha^k.
\end{equation}
The right side denotes the integral of any smooth representative of
\(\alpha\); Stokes' theorem makes it independent of that choice.
Neither subspace needs to be normal.

If \(V'\not\subset\Supp P_i\), take the projective resolution
\(r:U\to V'\) with smooth K\"ahler source used above.  The pullback of
\(s\) is nonzero and defines
the effective rational divisor \(E=m^{-1}\operatorname{div}(r^*s)\)
on \(U\).  Put \(a=r^*(\alpha|_{V'})\) and
\(w=r^*[\omega_i|_{V'}]\).  The class \(a\) is nef, \(w\) has a
semipositive smooth representative, and the actual section gives
\(a-w=c_1(E)\).  Therefore
\begin{equation}\label{eq:outer-outside-support}
 \int_U a^k-\int_U w^k
 =\sum_{j=0}^{k-1}c_1(E)\cdot a^{k-1-j}w^j\geq0.
\end{equation}
For each summand, choose a smooth representative of
\(a+\varepsilon[\eta]\) which is semipositive, where \(\eta\) is a
fixed K\"ahler form on \(U\).  Restrict it and the semipositive
representative of \(w\) to every effective divisor component, integrate,
and let \(\varepsilon\) tend to zero.
This proves its nonnegativity, including the degree statement when
\(k=1\).  Projection formula gives
\(\int_U w^k=\int_{V'_{\rm reg}}\omega_i^k>0\).  Equations
\eqref{eq:outer-strict-transform}--\eqref{eq:outer-outside-support}
give the desired strict positivity.

It remains that \(V'\) lies in a component \(T'\) of \(\Supp P_i\).
It meets \(X\setminus T\), so \(T'\neq T\).  The actual floor
specialization in Subsection~\ref{subsec:actual-floor-split} gives
\begin{equation}\label{eq:outer-other-floor}
 g:T'\to\overline W,\qquad \alpha|_{T'}=g^*\kappa,
\end{equation}
where \(g\) is projective with connected fibers, its source and target
are normal compact K\"ahler, and \(\kappa\) is a K\"ahler class.
The equality is in the local-potential Bott--Chern group.
Every irreducible curve in a \(g\)-fiber has \(\alpha\)-degree zero,
so \eqref{eq:global-support-zero-ray} puts its class in \(R_i\), and
\(c\) contracts it.

It follows that \(h=c|_{T'}\) is pointwise constant on every full
\(g\)-fiber.  To see this also for nonreduced fibers, their reductions
are connected projective varieties, possibly reducible.  If the image
of such a reduction were
not a point, one irreducible component \(Q\) would have nonconstant
image; otherwise its image would be a connected finite set.  At a smooth
point of \(Q\) where a local coordinate of \(h\) has nonzero
differential, general hyperplanes through that point cut an irreducible
projective curve component whose tangent is not in the differential's
kernel.
That curve has nonconstant image, a contradiction.  If \(\dim Q=1\),
take \(Q\) itself.

There is consequently a holomorphic factorization
\begin{equation}\label{eq:outer-factorization}
 h=b\circ g,\qquad b:\overline W\to Z.
\end{equation}
For clarity, the specialization gives
\(g_*\OO_{T'}=\OO_{\overline W}\); this also follows from analytic
Stein factorization, connected fibers, and the normality of the two
spaces.  Proper surjectivity makes \(g\) a closed quotient map, so the
pointwise constancy first defines a continuous \(b\).  For an open
\(O\subset Z\) embedded in a polydisc, restrict to the open
\(b^{-1}(O)\).  The coordinate functions of \(h\) on its full inverse
image descend through \(g_*\OO_{T'}=\OO_{\overline W}\) to a
holomorphic tuple.  It lies pointwise in the embedded \(O\);
its defining ideal therefore vanishes on the reduced open \(b^{-1}(O)\).
The resulting local holomorphic maps glue to
\eqref{eq:outer-factorization}.  This factorization uses pointwise
constancy on the underlying full fibers.

Put \(H=g(V')\), a closed irreducible analytic subspace of
\(\overline W\).  Since \(c|_{V'}\) is bimeromorphic onto \(V\),
\eqref{eq:outer-factorization} yields
\[
 k=\dim c(V')=\dim b(H)\leq\dim H\leq\dim V'=k.
\]
Thus \(g|_{V'}\) is generically finite onto \(H\), of some integer
degree \(d>0\).  The projection formula using
\eqref{eq:outer-other-floor} now gives
\begin{equation}\label{eq:outer-positive-intersections}
 \int_{V_{\rm reg}}\beta_0^k
 =\int_{V'_{\rm reg}}(g^*\kappa)^k
 =d\int_{H_{\rm reg}}\kappa^k>0.
\end{equation}
The last integral is a positive K\"ahler volume.  The one-way
factorization \(h=b\circ g\) and this degree calculation give the needed
implication in \cite[proof of Theorem~7.2, p.~48]{DasHaconPaun2024}.

These cases cover every irreducible positive-dimensional \(V\).
Additivity over top-dimensional irreducible components gives strict
positivity for every positive-dimensional compact reduced subspace.
We have proved analytic nefness of \(\beta\), the dominating current
\eqref{eq:outer-current-big}, and all positive top intersections.
The normal compact \(Z\) and the smooth real closed \(\beta_0\)
therefore satisfy the three hypotheses of
\cite[Theorem~2.29]{DasHaconPaun2024}.  That theorem gives a positive
smooth representative of \(\beta\).  Its local smooth potentials make
this a K\"ahler form \(\omega_{Z_i}\).  Equation
\eqref{eq:outer-descended-class} gives
\eqref{eq:outer-target-kahler}, and \(Z_i\) is compact K\"ahler.

\end{proof}

For a small step, the construction of its flipped morphism also has to
respect the hypotheses of the relative canonical-model theorem.
Use the current representative in
\eqref{eq:running-supported-representative}, and write
\(P_i=\sum p_jS_j\) with every \(p_j>0\) and \(S_j\) the floor primes.
At a nontrivial step this list is nonempty, since otherwise
\(A_i\sim_\Q0\) is nef.  Independently choose a rational
\(0<\varepsilon<1\) small
enough that
\[
 B_i^\varepsilon=B_i-\varepsilon P_i
\]
is effective.  Its floor coefficients are below one.  On a dlt resolution,
subtracting the pullback of the effective \(\Q\)-Cartier divisor
\(\varepsilon P_i\) can only lower the exceptional crepant coefficients,
which were already below one.  Hence \((X_i,B_i^\varepsilon)\) is klt.
Its actual adjoint satisfies
\begin{equation}\label{eq:flip-perturbation}
 K_{X_i}+B_i^\varepsilon=A_i-\varepsilon P_i
 \sim_\Q(1-\varepsilon)A_i.
\end{equation}
The projective contraction is bimeromorphic, so this adjoint is relatively
big.  Corollary~3.7 of \cite{DasHaconPaun2024} therefore applies at its
stated klt scope.  The actual isomorphism in \eqref{eq:flip-perturbation}
identifies common Cartier Veroneses of its relative algebra and of the
relative \(A_i\)-algebra.  On the compact base the latter is thus finitely
generated.  A further Veronese is generated in degree one, and its coherent
degree-one piece embeds its relative \(\operatorname{Projan}\) in the
associated relative projective space.  The tautological line is globally
relatively ample.  The ordinary small flip supplied by Theorem~7.2 is
this relative canonical model; in a common degree the tautological line
agrees with the flipped adjoint line on the common codimension-one open
and hence everywhere by reflexive extension.  The flipped morphism is
projective and the flipped adjoint is positive on its contracted curves.
Thus Corollary~3.7 is applied to the klt perturbation, whose common
Veronese is identified with that of the original adjoint.

The transform identities follow by induction from \(P_Y\), one completed
step at a time.  Divisorial contractions
remove their contracted primes and flips change no prime valuation.  On
the open containing all codimension-one points of the destination, the
fixed meromorphic pluriadjoint section identifies its line with the
strict transform \(P_{i+1}\).  After clearing the new indices, reflexive
extension gives
\[
 P_{i+1}\sim_\Q A_{i+1},\qquad
 \Supp P_{i+1}=\Supp\floor{B_{i+1}}.
\]
This establishes the identity needed for the perturbation at the next
step.  It uses ordinary \(\Q\)-factoriality only for the finitely many
global prime transforms.

We have now constructed a projective step, proved that its target is
K\"ahler, and supplied the effective actual representative needed to
repeat the construction.  The factorial/dlt preservation and the
special-termination argument in
\cite[proof of Theorem~7.2]{DasHaconPaun2024} apply to the unchanged
original boundary.  That construction, with
Proposition~\ref{prop:prime-floor-extension} at its ambient extension
steps and the target positivity proved above, therefore gives a finite
supported sequence to an analytically nef adjoint.

For a flip, take the main component of
\(X_i\times_{Z_i}X_{i+1}\) as its graph; its projections are projective.
For a divisorial contraction the source is its graph.  The main component
of the iterated fiber product of these finitely many graphs, followed by
normalization, is a normal common graph \(\Gamma\).  Normalization is
finite and projective, and projective morphisms compose over a compact
base \cite[Remark~2.11]{DasHaconPaun2024}.  Thus \(\Gamma\) is projective
over every running model.  Projective resolutions of it supply all common
resolutions with smooth source used below.  Their sources are compact K\"ahler by
the relative metric construction.

On a common projective resolution of one step with smooth source \(W_i\),
let \(r:W_i\to X_i\) and \(s:W_i\to X_{i+1}\) be the maps, and put
\begin{equation}\label{eq:step-difference}
 E_i=r^*P_i-s^*P_{i+1}.
\end{equation}
It is \(s\)-exceptional by the codimension-one comparison.  For an
\(s\)-contracted curve, its image under \(r\) is a point or a curve in
the negative source contraction fiber, so \(E_i\) has nonpositive degree.
Negativity gives \(E_i\geq0\).  In any common Cartier degree,
\begin{align*}
 H^0(X_i,\OO_{X_i}(mP_i))
 &=H^0(W_i,\OO_{W_i}(mr^*P_i))\\
 &=H^0(W_i,\OO_{W_i}(ms^*P_{i+1}+mE_i))\\
 &=H^0(X_{i+1},\OO_{X_{i+1}}(mP_{i+1})).
\end{align*}
The last equality is exceptional Hartogs extension for \(s\).
The identifications agree with the fixed section on the common open and
therefore respect multiplication.  A common Veronese section ring is
unchanged through the finite chain.  The power argument in
Lemma~\ref{lem:supported-preparation} handles other degrees, so the final
pair \((V,B)\), with \(A=K_V+B\) and \(P\) the final transform, satisfies
\begin{equation}\label{eq:final-supported}
 A\sim_\Q P\geq0,\qquad
 \Supp P=\Supp\floor B,\qquad \kappa(V,A)=0.
\end{equation}

\subsection{Discrepancy increase and a resolution of the floor}

We now show that the final pair has a resolution which is generically
unchanged along every floor intersection.  The point is that a negative
step strictly increases discrepancies over all its nontrivial fibers,
including those on the positive side of a flip.

\begin{lemma}[Support on a projective fiber]\label{lem:fiber-support}
Let \(h:W\to Z\) be a projective morphism between normal analytic spaces
with connected fibers, let
\(F\) be a scheme fiber, and let \(E\geq0\) be a rational
\(\Q\)-Cartier divisor on \(W\).  Suppose that \(E\cdot C\leq0\) for
every irreducible curve \(C\) in \(F\).  Then either
\(\Supp E\cap F=\varnothing\) or \(F\subseteq\Supp E\).
\end{lemma}
\begin{proof}
Clear the index and pass to the reduction of \(F\).  If an irreducible
component \(F_\alpha\) is not contained in \(\Supp E\) but meets it,
the restricted canonical section cuts a nonempty effective Cartier
divisor on the integral projective \(F_\alpha\).  This component has
positive dimension: a point that is an irreducible component cannot meet
another component, and a connected zero-dimensional fiber is a point.
Cutting by sufficiently general hyperplanes of a very ample line on
\(F_\alpha\) gives a curve on which \(E\) has strictly positive degree,
a contradiction.  If some components are contained in the support and
others are not, connectedness gives a noncontained component meeting a
contained one and the same contradiction.  These alternatives prove the
claim, with no reducedness assumption on the original fiber.
\end{proof}

Apply the lemma to a common resolution over the contraction base of one
step and to \(E_i\) in \eqref{eq:step-difference}.  The maps to the
contraction base have connected fibers: they are proper bimeromorphic
over a normal space.  On a base-fiber curve \(C\), a divisorial step has
\(E_i\cdot C=(r^*A_i)\cdot C\leq0\).  For a flip,
\[
 E_i\cdot C=(r^*A_i)\cdot C-(s^*A_{i+1})\cdot C\leq0,
\]
because the source adjoint is negative on its contracted curves and the
flipped adjoint is positive on its contracted curves.  Above any point
where either side has a nontrivial fiber, a curve in the common fiber can
be chosen to dominate a curve of that side: take a component over the
curve and cut by relative ample hyperplanes.  The corresponding term in
the preceding degree is then strictly negative.  Thus \(E_i\) meets the
common fiber, and Lemma~\ref{lem:fiber-support} puts the entire fiber in
\(\Supp E_i\).

Use the fixed meromorphic pluriadjoint section to calculate discrepancies
on the resolution.  Its divisor as a section of the old pulled-back line
is \(r^*P_i\), and as a section of the new line is \(s^*P_{i+1}\).
The two meromorphic pluriforms agree.  Their difference is consequently
the difference of the two crepant boundary divisors, so for every prime
\(F\)
\begin{equation}\label{eq:discrepancy-increase}
 \operatorname{coeff}_F(E_i)
 =a(F;X_{i+1},B_{i+1})-a(F;X_i,B_i).
\end{equation}
The same formula applies on higher projective resolutions after pullback.

For each contraction base \(Z_i\), let \(C_i\subset Z_i\) consist of
the points whose fiber on the source or, for a flip, on the destination
has positive dimension.  This is a closed analytic subset of the base
by properness and the fiber-dimension theorem.  A proper bimeromorphic
morphism to a normal space is an
isomorphism near any zero-dimensional fiber, since it is finite there.
Write \(v:\Gamma\to V\) and \(q_i:\Gamma\to Z_i\) for the projections
from the normal common graph to the final model and the contraction bases.
Set
\begin{equation}\label{eq:graph-bad-set}
 D_\Gamma=\bigcup_i q_i^{-1}(C_i),\qquad
 B_V=v(D_\Gamma),\qquad U_V=V\setminus B_V.
\end{equation}
The sets \(D_\Gamma\) and \(B_V\) are closed analytic, the latter by
properness.  We claim that
\begin{equation}\label{eq:graph-saturation}
 v^{-1}(B_V)=D_\Gamma.
\end{equation}
Indeed, at a point \(\gamma\notin D_\Gamma\), both sides of every step
are isomorphisms near the corresponding point of \(Z_i\).  Restricting
to these neighborhoods makes their iterated graph a common normal open,
so normalization does not change it and \(v\) is an isomorphism near
\(\gamma\).  Thus \(\{\gamma\}\) is open in its \(v\)-fiber, and it is
closed because the spaces are Hausdorff.  The fibers of \(v\) are
connected: a proper bimeromorphic morphism between normal spaces has
\(v_*\OO_\Gamma=\OO_V\), and its Stein factorization has connected
fibers.  That fiber is therefore \(\{\gamma\}\), so it cannot meet
\(D_\Gamma\).  This proves \eqref{eq:graph-saturation}.

The corresponding open \(U_Y\subset Y\) is isomorphic to \(U_V\)
through every step, because every successive lift there is unique.
No step extracts a divisor, so the
generic point of any prime on \(V\) follows a surviving prime through
the chain and avoids the contraction centers.  Therefore
\begin{equation}\label{eq:bad-codim}
 \operatorname{codim}_V(V\setminus U_V)\geq2.
\end{equation}

Every final log-canonical center tested by a prime on a projective
resolution meets \(U_V\).  To see this, suppose the final discrepancy
of that prime is zero.  Initial log canonicity, \(E_i\geq0\), and
\eqref{eq:discrepancy-increase} force all intermediate discrepancies
and all coefficients in the \(E_i\) to be zero.  If its center lay over
some \(C_i\), the fiber-support conclusion would place that center in
\(\Supp E_i\).  A local equation of a positive Cartier multiple of
\(E_i\) would then have positive order in the valuation, a contradiction.
If its center on \(V\) were contained in \(B_V\), its center on the
proper graph \(\Gamma\) would be contained in
\(v^{-1}(B_V)=D_\Gamma\).  Irreducibility would then put that entire
center in one \(q_i^{-1}(C_i)\).  On a higher projective resolution
carrying the prime, the center would lie over \(C_i\), giving the
contradiction just proved.  Thus it meets
\(U_V\), as claimed.  We only need this statement for primes on the
projective resolutions constructed here and for the blowups of their
floor strata.

\begin{lemma}[A globally simple dlt resolution]\label{lem:special-resolution}
The final pair \((V,B)\) has a projective log resolution
\(\pi:\widehat V\to V\) with the following properties:
\begin{enumerate}
 \item the strict boundary and exceptional support have globally smooth
 distinct SNC components;
 \item every \(\pi\)-exceptional crepant coefficient is strictly below one;
 \item if \(\widehat Z\) is an irreducible component of the nonempty
 intersection of distinct strict transforms of floor primes, then
 \(\pi\) is generically an isomorphism on \(\pi(\widehat Z)\).
\end{enumerate}
Moreover \(\pi\) is an isomorphism over \(U_V\).
\end{lemma}
\begin{proof}
Write \(b:\Gamma\to Y\) for the normal common graph.  Choose a relatively
very ample and generated line \(\mathscr L\) for \(b\); compactness permits
one global power.  Its evaluation embeds \(\Gamma\) in
\(\PP_Y(\mathscr E)\), where \(\mathscr E=b_*\mathscr L\).  This coherent
sheaf is torsion-free of rank one: a section killed by a nonzero function
vanishes on the dense isomorphism locus and hence everywhere.  Since \(Y\)
is smooth, \(\mathscr E^{**}\) is a line, and
\[
 \mathscr J=\mathscr E\otimes(\mathscr E^{**})^{-1}\subset\OO_Y
\]
is a coherent ideal equal to \(\OO_Y\) on \(U_Y\).

Apply the ideal principalization functor of
\cite[Theorems~1.1.11 and~1.1.13]{Temkin2017} to \(\mathscr J\).
For smooth input it is supported on the cosupport of the ideal, so it
avoids \(U_Y\); its total transform is invertible on a smooth space
\(Y_1\).  The resulting invertible quotient of the pulled-back
\(\mathscr E\) defines a map \(Y_1\to\PP_Y(\mathscr E)\).  It lands
in the closed subspace \(\Gamma\) on a dense open, and hence everywhere
because \(Y_1\) is reduced.  The map \(Y_1\to\Gamma\) is projective:
its graph is a closed immersion into the base change of the projective
map \(Y_1\to Y\).  Its composite to \(V\) is projective and is an
isomorphism over \(U_V\).

This principalization is not required to have transverse intermediate
centers.  Instead, after it is complete, mark as separate ordered Cartier
boundary components all strict initial boundary primes, all primes in the
principalized support, and all divisorial exceptional primes.  On \(U_Y\)
this is the original globally simple SNC list.  Apply the ordered-boundary
functor of \cite[Theorems~1.1.3 and~1.1.13]{Temkin2017}.  Its centers
lie over the original labeled bad locus and therefore avoid \(U_Y\).
The final indexed components and intersections are smooth, and the total
support consists of the strict support and new exceptional support.
The already invertible ideal remains invertible, so the map to \(\Gamma\)
persists.

Finally apply the same two stages to the pullback of the reduced coherent
ideal of \(V\setminus U_V\).  In the second stage mark, each distinct
prime once, every component of the prior resolved support, every component
of the newly principalized support, and every divisorial exceptional
prime.  Its total transform has divisorial support equal to
the entire inverse image of \(V\setminus U_V\); every such divisor maps
to a subset of codimension at least two by \eqref{eq:bad-codim}.  Outside
this support the resulting map \(\pi:\widehat V\to V\) is an
isomorphism.  Conversely a point on this support cannot be a local
isomorphism point, because the divisor germ through it would map to a
hypersurface contained in \(V\setminus U_V\).  Thus this is exactly the
nonisomorphism locus.  The complete support is globally simple SNC and
includes every strict boundary component of \(V\).

Each \(\pi\)-exceptional prime has center in \(V\setminus U_V\).  The
preceding discrepancy argument shows that its log discrepancy is positive,
proving property (2).  Let \(\widehat Z\) be an irreducible intersection
of \(k\) distinct transformed floor components.  At its general point
exactly those components occur, by SNC.  If \(k=1\), its prime valuation
has log discrepancy zero.  If \(k\geq2\), the blowup of the smooth stratum
has log discrepancy \(\operatorname{codim}(\widehat Z)-k=k-k=0\).
This is a prime on a projective resolution, so its image meets \(U_V\).
Since \(\pi\) is an isomorphism there, it is generically an isomorphism
on that image.  This proves property (3).
\end{proof}

\subsection{The support survives}

\begin{proof}[Proof of Proposition~\ref{prop:supported-model}]
All statements except \(P\ne0\) have been established.  Suppose that
\(P=0\).  Every component of \(P_Y\), and hence every component of
\(p^*M\), has disappeared.  On a common projective resolution with smooth source
\(r:W\to Y\), \(s:W\to V\), set \(Q=r^*p^*M\).  It is effective and
nonzero: the pullback of a nonzero effective Cartier multiple has a
nonzero strict transform.  It is \(s\)-exceptional.  Indeed, a prime of
\(W\) mapping onto a prime of \(V\) corresponds, by no extraction, to
a surviving prime of \(Y\); its coefficient in \(Q\) is zero under the
supposition that every component of \(p^*M\) disappeared.

The actual rational line of \(Q\) is \((pr)^*D\).  It is analytically
nef on the compact K\"ahler \(W\), by pullback of the metric inequalities
for the original \(D\).  In particular it is relatively nef for the
projective \(s\).  Exceptional negativity gives \(Q\leq0\), contradicting
its effectivity and nonvanishing.  Thus \(P\ne0\).
\end{proof}

The construction uses the original section but supplies no new ambient
section.  Its floor line is already defined by restricting a Cartier
multiple of \(A\) to the reduced subspace \(S=\floor B\).  The next
sections will generate that line on all of \(S\), after which supported
lifting contradicts the last equality in
Proposition~\ref{prop:supported-model}.
